\chapter{How to read this book}
The intended reader begins with ordinary English and curiosity about neural signals. Part I requires no mathematical notation beyond what it introduces. Parts II and III build number systems, limits, and finite-dimensional algebra before the book assumes calculus, probability, or complex phase. Later chapters use the same format: meaning in words, formal definition, calculation or derivation, a neuroscience connection, a failure mode, and exercises. A novice may work every exercise; an experienced reader can use the cross-reference path below to fill only missing prerequisites.

\section*{Suggested prerequisite chains}
\begin{description}[leftmargin=2cm,style=nextline]
\item[Hilbert and PAC] English $\to$ functions $\to$ complex numbers $\to$ sequences/integrals $\to$ Fourier $\to$ analytic signal $\to$ circular inference.
\item[PCA and manifolds] Relations $\to$ functions $\to$ real numbers $\to$ vector spaces $\to$ covariance $\to$ SVD/PCA $\to$ dynamical systems.
\item[Spike GLMs] Sets $\to$ functions $\to$ counting $\to$ probability $\to$ conditional probability $\to$ Poisson processes $\to$ likelihood and GLM.
\item[Granger analysis] Covariance $\to$ stochastic processes $\to$ cross spectra $\to$ multivariate linear prediction $\to$ null models and validation.
\item[Current source density] Coordinates $\to$ derivatives $\to$ multivariable calculus $\to$ forward measurement physics $\to$ regularized inverse methods.
\end{description}

\section*{Technical vocabulary}
An \emph{estimand} is the quantity the analysis intends to estimate; an \emph{estimator} is a data-dependent rule for estimating it. A \emph{model} is a set of assumed mathematical relationships; a \emph{measurement} is an observation through an instrument. \emph{Association} is not automatically \emph{causation}. \emph{Prediction} need not identify an intervention effect. The word ``significant'' should be accompanied by the hypothesis and statistical criterion being used.

\section*{About the two foundational books}
The structure draws inspiration from Halmos' conversational construction of sets and relations \cite{HalmosNaive}, and Mendelson's progression from logic and set theory to the number systems and real-number completeness \cite{MendelsonNumbers}. The explanations, examples, and proofs here are written afresh for CON$^2$EPHYS and are not reproduced sections of either text.